% GENERATED FILE. Edit canonical lessons, then run npm run generate:books.
% canonical-source-hash: fnv1a64:b25274e822f2dbdf
% canonical-lessons: ZH-C185-haojiu, ZH-C185-zhiqian, ZH-C185-zhihou, ZH-C185-zaowan, ZH-C185-wangwang

\chapter{Long Ago and Soon}
\label{ch:zh-long-ago-and-soon}
\hlchaptermodality{185}

\begin{chapteropening}
\textbf{By the end of this chapter:} \emph{I can say it has been ages, talk about before and after, mornings and evenings, and what tends to happen.}

Complete the last lesson of chapter 185: I can say it has been ages, talk about before and after, mornings and evenings, and what tends to happen.
\end{chapteropening}

\section[hǎojiǔ]{hǎojiǔ — a long time; ages}

\label{lesson:ZH-C185-haojiu}



\begin{quote}
Before the new one: say the Mandarin for one's home town, then the Mandarin for to relax.
\end{quote}

\subsection*{You'll want to know: hǎojiǔ}
\emph{hǎojiǔ} — \textquotedblleft{}a long time; ages\textquotedblright{}.

\begin{quote}
\emph{wǒmen hǎojiǔ méi jiànmiàn le} — We haven't seen each other for ages.
\end{quote}

Its characters come later; for now it is a word for the ear.

Two third tones: hǎojiǔ is said háojiǔ, the first rising.

\begin{cousinweb}
\emph{hǎojiǔ} is built from \zh{好} 'good', here 'very', and \emph{jiǔ} 'long-lasting', both inherited from Middle Chinese, \emph{xawX} and \emph{kjuwX} (Guangyun, in Baxter's transcription), which glosses the second 'lasting long'.
\end{cousinweb}

\begin{grammarlens}[title={What you've built}]
One more word for times of day.
\end{grammarlens}

\subsection*{Your turn}
\begin{itemize}
\raggedright
  \item \emph{Say it:} \emph{hǎojiǔ}
  \item \emph{Say it:} \emph{hǎojiǔ}, once more
  \item \emph{Say it:} \emph{wǒmen hǎojiǔ méi jiànmiàn le} — We haven't seen each other for ages.
  \item \emph{Recall:} say the Mandarin for one's home town, then the Mandarin for to relax, then say \emph{hǎojiǔ} again
\end{itemize}

\begin{tcolorbox}[breakable,colback=teal!4,colframe=teal!35!black,title={Before you move on}]
What does hǎojiǔ mean? (\textquotedblleft{}A long time; ages\textquotedblright{}.) Say it once more.
\end{tcolorbox}
\section[zhīqián]{zhīqián — before (a time); before (an event)}

\label{lesson:ZH-C185-zhiqian}



\begin{quote}
Before the new one: say the Mandarin for to relax, then the Mandarin for for ages.
\end{quote}

\subsection*{You'll want to know: zhīqián}
\emph{zhīqián} — \textquotedblleft{}before (a time); before (an event)\textquotedblright{}.

\begin{quote}
\emph{chī fàn zhīqián yào xǐ shǒu} — Wash your hands before you eat.
\end{quote}

Its characters come later; for now it is a word for the ear.

\begin{cousinweb}
\emph{zhīqián} is built from \emph{zhī}, the old word that links a thing to what it belongs to, and \zh{前} 'before', both inherited from Middle Chinese, \emph{tsyi} and \emph{dzen} (Guangyun, in Baxter's transcription). The dictionary glosses the first 'to go to'; the linking sense is not given there.
\end{cousinweb}

\begin{grammarlens}[title={What you've built}]
One more word for times of day.
\end{grammarlens}

\subsection*{Your turn}
\begin{itemize}
\raggedright
  \item \emph{Say it:} \emph{zhīqián}
  \item \emph{Say it:} \emph{zhīqián}, once more
  \item \emph{Say it:} \emph{chī fàn zhīqián yào xǐ shǒu} — Wash your hands before you eat.
  \item \emph{Recall:} say the Mandarin for to relax, then the Mandarin for for ages, then say \emph{zhīqián} again
\end{itemize}

\begin{tcolorbox}[breakable,colback=teal!4,colframe=teal!35!black,title={Before you move on}]
What does zhīqián mean? (\textquotedblleft{}Before (a time); before (an event)\textquotedblright{}.) Say it once more.
\end{tcolorbox}
\section[zhīhòu]{zhīhòu — after (a time); after (an event)}

\label{lesson:ZH-C185-zhihou}



\begin{quote}
Before the new one: say the Mandarin for for ages, then the Mandarin for before (a time).
\end{quote}

\subsection*{You'll want to know: zhīhòu}
\emph{zhīhòu} — \textquotedblleft{}after (a time); after (an event)\textquotedblright{}.

\begin{quote}
\emph{xiàbān zhīhòu wǒ qù mǎi dōngxi} — After work I'm going shopping.
\end{quote}

Its characters come later; for now it is a word for the ear.

\begin{cousinweb}
\emph{zhīhòu} is built from \emph{zhī}, the old linking word, and \zh{后} 'after', both inherited from Middle Chinese, \emph{tsyi} and \emph{huwX} (Guangyun, in Baxter's transcription).
\end{cousinweb}

\begin{grammarlens}[title={What you've built}]
One more word for times of day.
\end{grammarlens}

\subsection*{Your turn}
\begin{itemize}
\raggedright
  \item \emph{Say it:} \emph{zhīhòu}
  \item \emph{Say it:} \emph{zhīhòu}, once more
  \item \emph{Say it:} \emph{xiàbān zhīhòu wǒ qù mǎi dōngxi} — After work I'm going shopping.
  \item \emph{Recall:} say the Mandarin for for ages, then the Mandarin for before (a time), then say \emph{zhīhòu} again
\end{itemize}

\begin{tcolorbox}[breakable,colback=teal!4,colframe=teal!35!black,title={Before you move on}]
What does zhīhòu mean? (\textquotedblleft{}After (a time); after (an event)\textquotedblright{}.) Say it once more.
\end{tcolorbox}
\section[\texorpdfstring{\zh{早晚} (zǎowǎn)}{早晚 (zǎowǎn)}]{\zh{早晚} (zǎowǎn) — morning and evening; sooner or later}

\label{lesson:ZH-C185-zaowan}



\begin{quote}
Before the new one: say the Mandarin for before (a time), then the Mandarin for after (a time).
\end{quote}

\subsection*{You'll want to know: \zh{早晚}}
\textbf{\zh{早晚}} — \emph{zǎowǎn} — \textquotedblleft{}morning and evening; sooner or later\textquotedblright{}.

\begin{quote}
\textbf{\zh{这儿早晚都很冷}} \emph{zhèr zǎowǎn dōu hěn lěng} — Here it's cold morning and evening.
\end{quote}

Two third tones: zǎowǎn is said záowǎn, the first rising.

\begin{cousinweb}
\textbf{\zh{早晚}} (\emph{zǎowǎn}) is built from \zh{早} 'early' and \zh{晚} 'late', both inherited from Middle Chinese, \emph{tsawX} and \emph{mjonX} (Guangyun, in Baxter's transcription): early and late, and so sooner or later.
\end{cousinweb}

\begin{grammarlens}[title={What you've built}]
One more word for times of day.
\end{grammarlens}

\subsection*{Your turn}
\begin{itemize}
\raggedright
  \item \emph{Say it:} \emph{zǎowǎn}
  \item \emph{Say it:} \emph{zǎowǎn}, once more
  \item \emph{Say it:} \emph{zhèr zǎowǎn dōu hěn lěng} — Here it's cold morning and evening.
  \item \emph{Recall:} say the Mandarin for before (a time), then the Mandarin for after (a time), then say \emph{zǎowǎn} again
\end{itemize}

\begin{tcolorbox}[breakable,colback=teal!4,colframe=teal!35!black,title={Before you move on}]
What does \zh{早晚} mean? (\textquotedblleft{}Morning and evening; sooner or later\textquotedblright{}.) Say it once more.
\end{tcolorbox}
\section[wǎngwǎng]{wǎngwǎng — more often than not; tends to}

\label{lesson:ZH-C185-wangwang}



\begin{quote}
Before the new one: say the Mandarin for after (a time), then the Mandarin for morning and evening.
\end{quote}

\subsection*{You'll want to know: wǎngwǎng}
\emph{wǎngwǎng} — \textquotedblleft{}more often than not; tends to\textquotedblright{}.

\begin{quote}
\emph{zhōumò wǒ wǎngwǎng shí diǎn qǐchuáng} — At weekends I tend to get up at ten.
\end{quote}

Its characters come later; for now it is a word for the ear.

Two third tones: wǎngwǎng is said wángwǎng, the first rising.

\begin{cousinweb}
\emph{wǎngwǎng} says one syllable twice. It is inherited from Middle Chinese \emph{hjwangX} (Guangyun, in Baxter's transcription), glossed 'to go'; how the doubled word came to mean 'more often than not' is not settled.
\end{cousinweb}

\begin{grammarlens}[title={What you've built}]
One more word for times of day.
\end{grammarlens}

\subsection*{Your turn}
\begin{itemize}
\raggedright
  \item \emph{Say it:} \emph{wǎngwǎng}
  \item \emph{Say it:} \emph{wǎngwǎng}, once more
  \item \emph{Say it:} \emph{zhōumò wǒ wǎngwǎng shí diǎn qǐchuáng} — At weekends I tend to get up at ten.
  \item \emph{Recall:} say the Mandarin for after (a time), then the Mandarin for morning and evening, then say \emph{wǎngwǎng} again
\end{itemize}

\begin{tcolorbox}[breakable,colback=teal!4,colframe=teal!35!black,title={Before you move on}]
What does wǎngwǎng mean? (\textquotedblleft{}More often than not; tends to\textquotedblright{}.) Say it once more.
\end{tcolorbox}
